This project already contains a working Python implementation in the file \texttt{src/gmm\_missing.py}, which makes it possible to connect the mathematical derivation to concrete engineering choices. That link is valuable because several numerically important details do not appear explicitly in the original paper.

\subsection{Algorithm Box}
\begin{figure}[t]
\hrule
\vspace{0.5em}
\begin{algorithmic}[1]
\Require incomplete matrix $\mat{X}^{\text{inc}}$, number of clusters $\scal{K}$, tolerance $\scal{\varepsilon}$, maximum iterations $\scal{T}_{\max}$
\Ensure completed matrix $\mat{X}$, responsibilities $\mat{\Gamma}$, labels $\vect{\lambda}$
\State initialize $\mat{X}$ by column-mean imputation
\State initialize $\set*{\scal{\alpha}_k,\vect{\mu}_k,\mat{\Sigma}_k}_{k=1}^{\scal{K}}$
\For{$\scal{t}=1,\dots,\scal{T}_{\max}$}
    \State compute responsibilities $\scal{\gamma}_{jk}$ for all $(j,k)$
    \State update $\scal{\alpha}_k$, $\vect{\mu}_k$, and $\mat{\Sigma}_k$
    \For{each sample $\vect{x}_j$ with missing block $\vect{x}_{j,m}$}
        \State update $\vect{x}_{j,m}$ by the precision-weighted closed form
        \State keep $\vect{x}_{j,o}$ unchanged
    \EndFor
    \If{relative log-likelihood improvement $\le \scal{\varepsilon}$}
    \State \textbf{break}
    \EndIf
\EndFor
\State $\lambda_j \gets \argmax_k\ \scal{\gamma}_{jk}$ for each sample
\end{algorithmic}
\vspace{0.5em}
\hrule
\caption{One-stage GMM clustering with incomplete data.}
\label{alg:gmm-missing}
\end{figure}

\subsection{Numerical Safeguards}
The code-level implementation introduces several safeguards that a professional report should document explicitly:
\begin{itemize}
    \item \textbf{Mean initialization for missing values.} A simple initial completion is used only to start the alternating routine; it is not the final imputation mechanism.
    \item \textbf{K-means parameter initialization.} Component means are initialized from K-means centers, which generally improves stability over random initialization.
    \item \textbf{Covariance regularization.} A diagonal ridge term is added to each covariance estimate to prevent singularity when some clusters are small or nearly collinear.
    \item \textbf{Log-domain responsibility computation.} Using log-likelihoods and a log-sum-exp normalization is more stable than evaluating Gaussian densities directly.
    \item \textbf{Fallback for singular precision blocks.} If a block inversion fails during the update of $\vect{x}_{j,m}$, the implementation retains the current value rather than crashing.
\end{itemize}

\subsection{Mapping the Equations to the Code}
The implementation mirrors the paper almost line by line. The E-step computes a matrix $\mat{\Gamma} \in \R^{\scal{n}\times \scal{K}}$ whose entries are the responsibilities $\scal{\gamma}_{jk}$. The parameter M-step then forms
\[
    \scal{N}_k = \sum_{j=1}^{\scal{n}} \scal{\gamma}_{jk},
    \qquad
    \vect{\mu}_k = \frac{1}{\scal{N}_k}\sum_{j=1}^{\scal{n}}\scal{\gamma}_{jk}\vect{x}_j,
\]
and
\[
    \mat{\Sigma}_k
    =
    \frac{1}{\scal{N}_k}
    \sum_{j=1}^{\scal{n}}
    \scal{\gamma}_{jk}
    \paren*{\vect{x}_j-\vect{\mu}_k}
    \paren*{\vect{x}_j-\vect{\mu}_k}\trans
    + \scal{\lambda}\mat{I},
\]
where $\scal{\lambda}>0$ is the regularization level used in code.

The data-update step then loops over samples, extracts missing and observed coordinate sets, partitions the component precision matrices accordingly, accumulates the weighted system matrix and right-hand side, and solves for $\vect{x}_{j,m}$. This is exactly the operational form of the paper's Equation (14).

\subsection{Reproducibility Recommendations}
To make the report more reproducible than the source paper, future experiments should log at least the following metadata:
\begin{itemize}
    \item random seed for missing-pattern generation and model initialization;
    \item covariance regularization level and stopping tolerance;
    \item average runtime and iteration count for each dataset and missing ratio;
    \item failure counts due to singular updates or degenerate clusters;
    \item sensitivity curves with respect to $\scal{K}$, initialization, and missingness severity.
\end{itemize}

\begin{remark}
If these metadata are omitted, a method may appear empirically strong while actually depending on a narrow band of stable hyperparameter settings. That risk is especially serious for non-convex mixture models.
\end{remark}
